\section{Exterior powers and the obstruction class}\label{sec:duality}

\subsection{Projective tensor and exterior powers}\label{ob:ssec:powers}

For $k\ge1$ give the algebraic exterior power $\Lambda^k_KE$ the seminorm
\[
 \norm z_\pi=\inf_{z=\sum_jx_{j1}\wedge\cdots\wedge x_{jk}}
                 \sum_j\prod_{a=1}^k\norm{x_{ja}}.
\]
Its quotient by the zero-seminorm subspace is the normed space $\Lambda^k_\pi E$, with universal alternating map
$\omega_E(x_1,\ldots,x_k)=[x_1\wedge\cdots\wedge x_k]$. Put $\Lambda^0_\pi E=K$ and $\omega_E()=1$.

\begin{proposition}[representation]\label{prop:exterior}
There is a natural linear isometry
\[
 \Alt^k(E;F)\cong\cL(\Lambda^k_\pi E,F)
\]
for every normed $F$. Under it, precomposition by $u\colon E\to D$ is precomposition by
$\Lambda^k_\pi u\colon\Lambda^k_\pi E\to\Lambda^k_\pi D$.
\end{proposition}
\begin{proof}
The algebraic linearization $\widetilde m$ of a bounded alternating map $m$ satisfies
$\norm{\widetilde m(z)}\le\norm m\norm z_\pi$, by the definition of $\norm z_\pi$; so it kills the null subspace and
descends contractively. Conversely, $\norm{\omega_E(x_1,\ldots,x_k)}\le\prod_a\norm{x_a}$, so composing a bounded linear
map with $\omega_E$ is contractive. The two operations are inverse, hence isometric. Applying $u$ to each
wedge factor gives the intertwining identity and $\norm{\Lambda^k_\pi u}\le\norm u^k$ for $k\ge1$; in degree zero the
action is the identity on $K$.
\end{proof}

If $F$ is complete, the same correspondence holds with the completion $\widehat{\Lambda^k_\pi E}$.

The projective tensor product $X_1\otimes_\pi\dots\otimes_\pi X_n$ of normed spaces, and $\otimes^k_\pi E$ with $k$
factors $E$, are defined in the same way with $\otimes$ in place of $\wedge$. As in Proposition~\ref{prop:exterior},
bounded $n$-linear maps $m$ on $X_1\times\dots\times X_n$ correspond isometrically to bounded linear maps $m^\otimes$ on
$X_1\otimes_\pi\dots\otimes_\pi X_n$. For bounded operators $u_i$, $u_1\otimes\dots\otimes u_n$ has norm at most
$\prod_i\norm{u_i}$ and is multilinear in $(u_1,\dots,u_n)$; $u^{\otimes k}=u\otimes\dots\otimes u$. Hats denote
completions, maps extend to completions by continuity, and every normed space is a subspace of its completion. Let
\[
 \wedge^k_E\colon\otimes^k_\pi E\to\Lambda^k_\pi E,\qquad
 \operatorname{alt}^k_E\colon\Lambda^k_\pi E\to\otimes^k_\pi E
\]
be the maps with $x_1\otimes\dots\otimes x_k\mapsto x_1\wedge\dots\wedge x_k$ and
$x_1\wedge\dots\wedge x_k\mapsto\sum_{\sigma\in S_k}\sgn(\sigma)\,x_{\sigma(1)}\otimes\dots\otimes x_{\sigma(k)}$. The first is
$\omega_E^\otimes$, a contraction. The second is well defined in every characteristic, since for $x_a=x_b$ the
permutations $\sigma$ and $(a\ b)\circ\sigma$ give equal tensors with opposite signs; its norm is at most the real number
$k!$, and $\wedge^k_E\circ\operatorname{alt}^k_E=k!$. Both maps are natural:
$\wedge^k_D\circ u^{\otimes k}=\Lambda^k_\pi u\circ\wedge^k_E$ and
$u^{\otimes k}\circ\operatorname{alt}^k_E=\operatorname{alt}^k_D\circ\Lambda^k_\pi u$ for $u\in\cL(E,D)$.

\subsection{Universal and scalar tests}\label{ob:ssec:tests}

Write $Y^*=\cL(Y,K)$ and $J_Y(y)(\lambda)=\lambda(y)$ for the canonical contraction into the continuous bidual; no
point-separation or Hahn--Banach assumption is implicit.

\begin{theorem}[universal and scalar tests]\label{thm:duality}
Let $U$ be an open parameter domain and $W\colon U\to\cL(X,Y)$ a family. Define
\[
 R_F(t)(T)=T\circ W(t),\qquad
 R_F\colon U\to\cL(\cL(Y,F),\cL(X,F)).
\]
For power-series analyticity, or for any fixed class $C^n$, the following statements hold.
\begin{enumerate}
\item $W$ is regular if and only if $R_F$ is regular for every normed $F$; the single target $Y$ suffices.
\item $R_K$ is regular if and only if $J_YW$ is regular in $\cL(X,Y^{**})$, and if and only if
$R_{G^*}$ is regular for every normed $G$. Any one nonzero continuous dual $G^*$ suffices for the last test.
\item If $J_Y$ has a bounded linear left inverse, regularity of $R_K$ implies regularity of $W$.
If $J_F$ has such a left inverse, regularity of $R_K$ implies regularity of $R_F$.
\end{enumerate}
\end{theorem}
\begin{proof}
All identities below use bounded linear operations, which preserve each regularity class. Precomposition depends
bounded-linearly on $W$, and $W(t)=R_Y(t)(\id_Y)$; this proves~(1). Exchanging the arguments defines a linear isometry
\[
 \Theta\colon\cL(Y^*,X^*)\longrightarrow\cL(X,Y^{**}),
 \qquad (\Theta(B)x)(\lambda)=(B\lambda)(x),
\]
since both operator norms are the least $C$ bounding the displayed scalar by $C\norm x\norm\lambda$; its inverse exchanges
the arguments again, and $\Theta(R_K(t))=J_YW(t)$. The isometry
$\cL(X,G^*)\cong\cL(G,X^*)$, $T\mapsto[g\mapsto(x\mapsto T(x)(g))]$, identifies $R_{G^*}(t)$ with $S\mapsto R_K(t)\circ S$,
which gives regularity for every dual target. Conversely, choose $v\in G^*$ and $g\in G$ with $v(g)=1$; the injection
$\alpha\mapsto v\alpha$ and evaluation at $g$ compress $R_{G^*}$ to $R_K$. This proves~(2). For~(3), if $P_YJ_Y=\id_Y$,
then $W=P_Y\Theta(R_K)$; if $P_FJ_F=\id_F$, then $R_F(t)(T)=P_F\circ R_{F^{**}}(t)(J_F\circ T)$.
\end{proof}

\subsection{Good pairs and one test target}\label{ob:ssec:good}

A \emph{Banach space} is a complete normed $K$-space.

\begin{definition}\label{fib:def:good}
A pair $(E,D)$ of normed spaces is \emph{good in degree $k$} if $A^k_{E,D;F}$ is analytic for every normed $F$, and
\emph{good for Banach values in degree $k$} if this holds for every Banach space $F$.
\end{definition}

\begin{lemma}[one test target]\label{fib:lem:test}
The pair $(E,D)$ is good in degree $k$ if and only if $u\mapsto\Lambda^k_\pi u$ is analytic at some point of $\cL(E,D)$,
and if and only if there is a bounded $k$\nobreakdash-linear $Q\colon\cL(E,D)^k\to\cL(\Lambda^k_\pi E,\Lambda^k_\pi D)$
with $Q(u,\dots,u)=\Lambda^k_\pi u$. The same holds for Banach values, with $\cL(\Lambda^k_\pi E,\widehat\Lambda^k_\pi D)$
in place of $\cL(\Lambda^k_\pi E,\Lambda^k_\pi D)$, also as the space in which $u\mapsto\Lambda^k_\pi u$ is analytic.
\end{lemma}

\begin{proof}
Apply Theorem~\ref{thm:duality}(1) with $X=\Lambda^k_\pi E$, $W(u)=\Lambda^k_\pi u$, and $Y=\Lambda^k_\pi D$,
respectively $Y=\widehat\Lambda^k_\pi D$. By Proposition~\ref{prop:exterior} and its completed form, $R_F$ is
$A^k_{E,D;F}$ for every normed $F$, respectively every Banach $F$, and the target $Y$ suffices. Proposition~\ref{prop:homogeneous-lift}
applies to $W$, with the closed isometry $T\mapsto T\circ\omega_E$ of $\cL(\Lambda^k_\pi E,Y)$ onto
$\Alt^k(E;Y)\subseteq\Mult^k(E;Y)$ and the ambient map $(u_1,\dots,u_k)\mapsto\omega_D\circ(u_1,\dots,u_k)$.
\end{proof}

We use two splitting conditions on normed spaces:
\begin{enumerate}
\item[$(\mathrm S^k_E)$] there is a bounded linear $s\colon\Lambda^k_\pi E\to\widehat\otimes^k_\pi E$ with
$\widehat\wedge^k_E\circ s=\id$;
\item[$(\mathrm R^k_D)$] there is a bounded linear $\rho\colon\otimes^k_\pi D\to\widehat\Lambda^k_\pi D$ with
$\rho\circ\operatorname{alt}^k_D=\id$.
\end{enumerate}
They hold \emph{without completion} if $s$ takes values in $\otimes^k_\pi E$, respectively $\rho$ in
$\Lambda^k_\pi D$. Both hold without completion if $k!\ne0$ in $K$, with $s=\operatorname{alt}^k_E/k!$ and
$\rho=\wedge^k_D/k!$, and if the space has finite continuous coordinates $(e_i,e^*_i)$, with
\[
 \begin{gathered}
 s(x_1\wedge\dots\wedge x_k)=\sum_{i_1<\dots<i_k}\det\bigl(e^*_{i_a}(x_b)\bigr)_{a,b}\,
 e_{i_1}\otimes\dots\otimes e_{i_k},\\
 \rho(y_1\otimes\dots\otimes y_k)=\sum_{i_1<\dots<i_k}\prod_a e^*_{i_a}(y_a)\,e_{i_1}\wedge\dots\wedge e_{i_k}.
 \end{gathered}
\]

\begin{proposition}[splittings]\label{fib:prop:split}
Let $E$ and $D$ be normed spaces.
\begin{enumerate}
\item If $(\mathrm S^k_E)$ holds, then $(E,D)$ is good for Banach values in degree $k$ for every normed $D$:
$Q(u_1,\dots,u_k)=\widehat\wedge^k_D\circ(u_1\otimes\dots\otimes u_k)\circ s$ is a lift as in
Lemma~\ref{fib:lem:test}, of norm at most $\norm s$.
\item If $(\mathrm R^k_D)$ holds, then $(E,D)$ is good for Banach values in degree $k$ for every normed $E$, with the
lift $Q(u_1,\dots,u_k)=\rho\circ(u_1\otimes\dots\otimes u_k)\circ\operatorname{alt}^k_E$ of norm at most $k!\norm\rho$.
\item If the condition holds without completion, the pairs in~(1), respectively~(2), are good in degree $k$.
\item $(\mathrm S^k_E)$ holds if and only if the inclusions $j_{E,F}\colon\Alt^k(E;F)\to\Mult^k(E;F)$, for all Banach
spaces $F$, have bounded linear retractions that commute with postcomposition by bounded linear maps $F\to F'$. It
holds without completion if and only if such retractions exist for all normed $F$.
\end{enumerate}
\end{proposition}

The proof is in Appendix~\ref{app:deferred}.

\subsection{Three exact sequences}\label{ob:ssec:sequences}

A sequence $\mathcal Y\colon0\to Y_1\xrightarrow{i}Y_2\xrightarrow{q}Y_3\to0$ of Banach spaces and bounded linear maps is
\emph{exact} if $i$ is injective, $i(Y_1)=\ker q$, and for a constant $C$, $\norm y\le C\norm{iy}$ on $Y_1$ and every
$z\in Y_3$ is $qy$ for some $y$ with $\norm y\le C\norm z$. It \emph{splits} if $q$ has a bounded linear right inverse.
Exact sequences $\mathcal Y,\mathcal Y''$ with the same end terms are \emph{equivalent} if a bounded linear
$T\colon Y_2\to Y''_2$ satisfies $Ti=i''$ and $q''T=q$; such a $T$ is bijective with bounded inverse, since $q$ is open and
$i''$ is bounded below, so this is an equivalence relation, and it preserves splitting. $\Ext(Y_3,Y_1)$ is the set of
classes $[\mathcal Y]$, and $0$ is the class of $0\to Y_1\to Y_1\oplus Y_3\to Y_3\to0$, with the inclusion, the projection
and the norm $\norm{y_1}+\norm{y_3}$. Then $[\mathcal Y]=0$ if and only if $\mathcal Y$ splits: a left inverse $\lambda$ of
$i$ (Lemma~\ref{ob:lem:split-criteria}(2)) gives the equivalence $y\mapsto(\lambda y,qy)$. For a bounded linear
$\alpha\colon Y_1\to Y'$ into a Banach space, the \emph{pushout} $\alpha_*\mathcal Y$ is
$0\to Y'\xrightarrow{i_\alpha}Y_\alpha\xrightarrow{q_\alpha}Y_3\to0$, where $Y_\alpha$ is the quotient of $Y'\oplus Y_2$,
normed by $\norm{y'}+\norm y$, by $\{(\alpha y,-iy):y\in Y_1\}$, which is closed because $y\mapsto(\alpha y,-iy)$ is
bounded below on the complete space $Y_1$; $i_\alpha y'=[(y',0)]$ and $q_\alpha[(y',y)]=qy$. Put
$\alpha_*[\mathcal Y]=[\alpha_*\mathcal Y]$; an equivalence $T$ induces the equivalence $[(y',y)]\mapsto[(y',Ty)]$, so
this depends only on $[\mathcal Y]$.

\begin{lemma}[splitting criteria]\label{ob:lem:split-criteria}
Let $\mathcal Y$ and $\mathcal Y'\colon0\to Y'_1\xrightarrow{i'}Y'_2\xrightarrow{q'}Y'_3\to0$ be exact.
\begin{enumerate}
\item Let $\alpha_a\colon Y_a\to Y'_a$, $a=1,2,3$, be bounded linear maps with $\alpha_2i=i'\alpha_1$ and
$q'\alpha_2=\alpha_3q$. Then $\alpha_1=\tilde\alpha_1i$ for a bounded linear $\tilde\alpha_1\colon Y_2\to Y'_1$ if and only
if $\alpha_3=q'\tilde\alpha_3$ for a bounded linear $\tilde\alpha_3\colon Y_3\to Y'_2$, and the two can be chosen with
$\alpha_2=i'\tilde\alpha_1+\tilde\alpha_3q$.
\item $\mathcal Y$ splits if and only if $i$ has a bounded linear left inverse.
\item For a bounded linear $\alpha\colon Y_1\to Y'$ into a Banach space, $\alpha_*\mathcal Y$ is exact, and
$\alpha_*[\mathcal Y]=0$ if and only if $\alpha$ extends to $Y_2$, that is, $\alpha=\tilde\alpha i$ for a bounded linear
$\tilde\alpha\colon Y_2\to Y'$.
\end{enumerate}
\end{lemma}

The proof is in Appendix~\ref{app:deferred}.

\begin{lemma}\label{ob:lem:metric-surjection}
For a normed space $E$ and $k\ge1$, $\wedge^k_E$ is a metric surjection: $\norm y=\inf\{\norm z:\wedge^k_Ez=y\}$ for
$y\in\Lambda^k_\pi E$. Its extension $\widehat\wedge^k_E\colon\widehat\otimes^k_\pi E\to\widehat\Lambda^k_\pi E$ is
surjective and again a metric surjection.
\end{lemma}

The proof is in Appendix~\ref{app:deferred}.

\begin{definition}\label{ob:def:sequences}
Let $E$, $D$, $P$ be normed spaces and $k,n\ge1$.
\begin{enumerate}
\item The \emph{source sequence} of $E$ in degree $k$ is the exact sequence (Lemma~\ref{ob:lem:metric-surjection})
\begin{equation}\label{ob:eq:source}
 \cS^k_E\colon\quad0\to\cN^k_E\to\widehat\otimes^k_\pi E\xrightarrow{\;\widehat\wedge^k_E\;}\widehat\Lambda^k_\pi E\to0,
 \qquad \cN^k_E=\ker\widehat\wedge^k_E,
\end{equation}
whose first map is the inclusion.
\item The \emph{target sequence} of $D$ in degree $k$ is
$\cR^k_D\colon0\to\widehat\Lambda^k_\pi D\xrightarrow{\alt^k_D}\widehat\otimes^k_\pi D\to
\widehat\otimes^k_\pi D/\overline{\alt^k_D(\widehat\Lambda^k_\pi D)}\to0$, with the quotient map last. It is exact if and
only if $\alt^k_D$ is bounded below on $\widehat\Lambda^k_\pi D$.
\item Let $\widehat\Gamma^n_\pi P$ be the closed linear span in $\widehat\otimes^n_\pi P$ of the pure powers
$h^{\otimes n}=h\otimes\dots\otimes h$, $h\in P$. The \emph{parameter sequence} of $P$ in degree $n$ is the exact sequence
$\cG^n_P\colon0\to\widehat\Gamma^n_\pi P\to\widehat\otimes^n_\pi P\to\widehat\otimes^n_\pi P/\widehat\Gamma^n_\pi P\to0$.
The space $P$ is \emph{polarizable in degree $n$}, written $(\Gamma^n_P)$, if $\cG^n_P$ splits, that is
(Lemma~\ref{ob:lem:split-criteria}(2)), if there is a bounded linear projection $R_n$ of $\widehat\otimes^n_\pi P$ onto
$\widehat\Gamma^n_\pi P$. Let $g_n(P)$ be the infimum of the norms of such projections.
\item For a Banach space $Y$, $\Mult^n_0(P;Y)$ is the closed subspace of $\Mult^n(P;Y)$ of the maps $\Phi$ with zero
diagonal, $\Phi(h,\dots,h)=0$ for $h\in P$, and $\Poly^n(P;Y)$ is the quotient of $\Mult^n(P;Y)$ by $\Mult^n_0(P;Y)$,
with the quotient norm; through $\Phi\mapsto(h\mapsto\Phi(h,\dots,h))$ its elements correspond to the diagonals of bounded
$n$-linear maps $P^n\to Y$, the $n$-homogeneous continuous polynomials of Section~\ref{sec:descent}. The resulting exact sequence is
\[
 \cM^n(P;Y)\colon\quad0\to\Mult^n_0(P;Y)\to\Mult^n(P;Y)\to\Poly^n(P;Y)\to0.
\]
\end{enumerate}
\end{definition}

\begin{lemma}\label{ob:lem:source}
Condition $(\mathrm S^k_E)$ holds if and only if $\cS^k_E$ in~\eqref{ob:eq:source} splits, and if and only if $\cN^k_E$ is complemented in
$\widehat\otimes^k_\pi E$. A section $s$ gives the projection $\id-s\,\widehat\wedge^k_E$ onto $\cN^k_E$, of norm at most
$1+\norm s$; a bounded projection $\Pi$ onto $\cN^k_E$ gives the section $s(\widehat\wedge^k_Ez)=(\id-\Pi)z$, of norm at
most $\norm{\id-\Pi}$.
\end{lemma}

\begin{proof}
A section on $\Lambda^k_\pi E$ extends by continuity to a right inverse of $\widehat\wedge^k_E$, and a right inverse
restricts to a section. The map $\id-s\,\widehat\wedge^k_E$ is the identity on $\cN^k_E$ and is killed by
$\widehat\wedge^k_E$. Conversely, $\id-\Pi$ vanishes on $\cN^k_E$, so the formula defines $s$, and
Lemma~\ref{ob:lem:metric-surjection} bounds its norm.
\end{proof}

\begin{lemma}\label{ob:lem:target}
For a normed space $D$ and $k\ge1$ the following are equivalent: \textup{(1)} $(\mathrm R^k_D)$; \textup{(2)} $\cR^k_D$ is
exact and splits; \textup{(3)} $\omega_D=\mathcal A_D\eta$ for some $\eta\in\Mult^k(D;\widehat\Lambda^k_\pi D)$, where
$(\mathcal A_D\eta)(y_1,\dots,y_k)=\sum_{\sigma\in S_k}\sgn(\sigma)\,\eta(y_{\sigma(1)},\dots,y_{\sigma(k)})$ is
unnormalized antisymmetrization.
\end{lemma}

The proof is in Appendix~\ref{app:deferred}.

\begin{lemma}\label{ob:lem:parameter}
For a normed space $P$ and $n\ge1$ the following are equivalent.
\begin{enumerate}
\item $P$ is polarizable in degree $n$.
\item The power map $P\to\widehat\Gamma^n_\pi P$, $h\mapsto h^{\otimes n}$, is the diagonal of a bounded $n$-linear map
$P^n\to\widehat\Gamma^n_\pi P$.
\item For every closed subspace $W$ of a Banach space $Z$, every bounded $n$-linear $\Phi\colon P^n\to Z$ with
$\Phi(h,\dots,h)\in W$ for all $h\in P$, and every $\delta>0$, there is a bounded $n$-linear $Q\colon P^n\to W$ with
$Q(h,\dots,h)=\Phi(h,\dots,h)$ for all $h$ and $\norm Q\le(g_n(P)+\delta)\norm\Phi$.
\end{enumerate}
Under these conditions $\cM^n(P;Y)$ splits for every Banach space $Y$, by right inverses that commute with postcomposition
by bounded linear maps $Y\to Y'$.
\end{lemma}

The proof is in Appendix~\ref{app:deferred}.

\begin{definition}\label{ob:def:split-operators}
Let $w\colon E\to G$ be a bounded linear map of normed spaces and $k\ge1$. Then $w$ is \emph{$\mathrm S^k$-split} if
$\widehat\wedge^k_G\circ s=\Lambda^k_\pi w$ for a bounded linear $s\colon\Lambda^k_\pi E\to\widehat\otimes^k_\pi G$, and
\emph{$\mathrm R^k$-split} if $\rho\circ\alt^k_E=\Lambda^k_\pi w$ for a bounded linear
$\rho\colon\otimes^k_\pi E\to\widehat\Lambda^k_\pi G$; it is split \emph{without completion} if $s$ takes values in
$\otimes^k_\pi G$, respectively $\rho$ in $\Lambda^k_\pi G$. It is \emph{$\mathrm S^{\le k}$-split}
(\emph{$\mathrm R^{\le k}$-split}) if it is $\mathrm S^j$-split ($\mathrm R^j$-split) for every $j\le k$.
\end{definition}

Thus $(\mathrm S^k_E)$ says that $\id_E$ is $\mathrm S^k$-split, and $(\mathrm R^k_D)$ that $\id_D$ is
$\mathrm R^k$-split. Every $w$ is $\mathrm S^1$- and $\mathrm R^1$-split, and $\mathrm S^j$- and $\mathrm R^j$-split
whenever $j!\ne0$ in $K$, with $s=w^{\otimes j}\alt^j_E/j!$ and $\rho=\wedge^j_Gw^{\otimes j}/j!$, since
$\wedge^j_Gw^{\otimes j}\alt^j_E=\Lambda^j_\pi w\,\wedge^j_E\alt^j_E=j!\,\Lambda^j_\pi w$.

\begin{proposition}[split operators]\label{ob:prop:split-operators}
Let $w\colon E\to G$ be a bounded linear map of normed spaces and $k\ge1$.
\begin{enumerate}
\item The map $w^{\otimes k}$ sends $\cN^k_E$ into $\cN^k_G$, and the following are equivalent: (a) $w$ is
$\mathrm S^k$-split; (b) $\widehat\Lambda^k_\pi w=\widehat\wedge^k_G\circ\tilde s$ for a bounded linear
$\tilde s\colon\widehat\Lambda^k_\pi E\to\widehat\otimes^k_\pi G$; (c) $w^{\otimes k}|_{\cN^k_E}\colon\cN^k_E\to\cN^k_G$
extends to a bounded linear map $\widehat\otimes^k_\pi E\to\cN^k_G$; (d) the \emph{splitting class}
$\operatorname{sp}_k(w)=(w^{\otimes k}|_{\cN^k_E})_*[\cS^k_E]\in\Ext(\widehat\Lambda^k_\pi E,\cN^k_G)$ vanishes. In
particular $(\mathrm S^k_E)$ holds if and only if $\operatorname{sp}_k(\id_E)=0$, if and only if $[\cS^k_E]=0$
(Lemma~\ref{ob:lem:source}).
\item If $w$ is $\mathrm S^k$-split with $s$, and $v\colon G\to G_1$, $v'\colon G_0\to E$ are bounded linear, then $vwv'$
is $\mathrm S^k$-split with $v^{\otimes k}s\,\Lambda^k_\pi v'$. If $w$ is $\mathrm R^k$-split with $\rho$, then $vwv'$ is
$\mathrm R^k$-split with $\widehat\Lambda^k_\pi v\,\rho\,v'^{\otimes k}$. The same holds without completion.
\end{enumerate}
\end{proposition}

The proof is in Appendix~\ref{app:deferred}.

\subsection{The obstruction class}\label{ob:ssec:class}

Throughout this subsection $E$ and $D$ are normed spaces and $k\ge1$; $\Mult^k(\cL(E,D);\widehat\Lambda^k_\pi D)$ is a
Banach space, and $\Mult^k_0(\cL(E,D);\widehat\Lambda^k_\pi D)$ is its closed subspace of maps with zero diagonal.

\begin{lemma}[evaluation]\label{ob:lem:evaluation}
There is a unique bounded linear map
$\ev=\ev^k_{E,D}\colon\widehat\otimes^k_\pi E\to\Mult^k(\cL(E,D);\widehat\Lambda^k_\pi D)$ with
\begin{equation}\label{ob:eq:evaluation}
 \ev(x_1\otimes\dots\otimes x_k)(u_1,\dots,u_k)=u_1x_1\wedge\dots\wedge u_kx_k,
\end{equation}
and $\norm\ev\le1$. It satisfies $\ev(z)(u,\dots,u)=\widehat\Lambda^k_\pi u(\widehat\wedge^k_Ez)$ for
$z\in\widehat\otimes^k_\pi E$ and $u\in\cL(E,D)$; in particular $\ev$ maps $\cN^k_E$ into
$\Mult^k_0(\cL(E,D);\widehat\Lambda^k_\pi D)$.
\end{lemma}

\begin{proof}
For $x\in E^k$ the right side of \eqref{ob:eq:evaluation} is $k$-linear in $(u_1,\dots,u_k)$, of norm at most
$\prod_a\norm{x_a}$, and $k$-linear in $x$. Its linearization in $x$ is a contraction from $\otimes^k_\pi E$ into a
Banach space and extends to $\widehat\otimes^k_\pi E$; pure tensors span a dense subspace, which gives uniqueness. For
fixed $u$, both sides of the identity are bounded linear in $z$ and agree on pure tensors.
\end{proof}

Let $q$ be the quotient map of $\cM^k(\cL(E,D);\widehat\Lambda^k_\pi D)$. Since $q\circ\ev$ vanishes on $\cN^k_E$,
Lemma~\ref{ob:lem:metric-surjection} gives a unique bounded linear
$\overline\ev\colon\hLambda^k_\pi E\to\Poly^k(\cL(E,D);\widehat\Lambda^k_\pi D)$ with
$q\circ\ev=\overline\ev\circ\widehat\wedge^k_E$. By Lemma~\ref{ob:lem:evaluation}, $\overline\ev(y)$ is the class of every
map in $\Mult^k(\cL(E,D);\widehat\Lambda^k_\pi D)$ with diagonal $u\mapsto\widehat\Lambda^k_\pi u(y)$.

The \emph{obstruction class} of the pair $(E,D)$ in degree $k$ is
\begin{equation}\label{ob:eq:class}
 \begin{split}
 \ob_k(E,D)&=(\ev|_{\cN^k_E})_*[\cS^k_E]\\
 &\in\Ext\bigl(\widehat\Lambda^k_\pi E,\Mult^k_0(\cL(E,D);\widehat\Lambda^k_\pi D)\bigr).
 \end{split}
\end{equation}

\begin{theorem}[obstruction theorem]\label{ob:thm:obstruction}
Let $E$ and $D$ be normed spaces and $k\ge1$. The following are equivalent.
\begin{enumerate}
\item The pair $(E,D)$ is good for Banach values in degree $k$.
\item The restriction $\ev|_{\cN^k_E}$ extends to a bounded linear map
$\widehat\otimes^k_\pi E\to\Mult^k_0(\cL(E,D);\widehat\Lambda^k_\pi D)$.
\item $\ob_k(E,D)=0$ in~\eqref{ob:eq:class}.
\end{enumerate}
More precisely, a lift $Q\colon\cL(E,D)^k\to\cL(\Lambda^k_\pi E,\widehat\Lambda^k_\pi D)$ as in
Lemma~\ref{fib:lem:test} gives an extension in~(2) of norm at most $1+\norm Q$, and an extension $\xi$ in~(2) gives such a
lift of norm at most $1+\norm\xi$.
\end{theorem}

\begin{proof}
By Lemma~\ref{fib:lem:test}, (1) holds if and only if a lift $Q$ exists.

\emph{Step 1 (lifts give extensions).} Bounded $k$-linear maps
$Q\colon\cL(E,D)^k\to\cL(\Lambda^k_\pi E,\widehat\Lambda^k_\pi D)$ correspond to bounded linear maps
$Q^\sharp\colon\widehat\Lambda^k_\pi E\to\Mult^k(\cL(E,D);\widehat\Lambda^k_\pi D)$, by
$Q^\sharp(y)(u_1,\dots,u_k)=Q(u_1,\dots,u_k)y$ for $y\in\Lambda^k_\pi E$, with $\norm{Q^\sharp}=\norm Q$. Given a lift $Q$,
put $\xi=\ev-Q^\sharp\widehat\wedge^k_E$; it equals $\ev$ on $\cN^k_E=\ker\widehat\wedge^k_E$. By
Lemma~\ref{ob:lem:evaluation}, $\xi(z)(u,\dots,u)=\Lambda^k_\pi u(\wedge^k_Ez)-Q(u,\dots,u)(\wedge^k_Ez)=0$ for
$z\in\otimes^k_\pi E$, hence for all $z$ by continuity. So $\xi$ is an extension as in~(2). Since $\norm\ev\le1$ by
Lemma~\ref{ob:lem:evaluation} and $\widehat\wedge^k_E$ is contractive, $\norm\xi\le\norm\ev+\norm{Q^\sharp}\le1+\norm Q$.

\emph{Step 2 (extensions give lifts).} Given an extension $\xi$ as in~(2), the map $\ev-\xi$ vanishes on $\cN^k_E$. Since $\widehat\wedge^k_E$ is a
metric surjection with kernel $\cN^k_E$ (Lemma~\ref{ob:lem:metric-surjection}), $\ev-\xi=Q^\sharp\widehat\wedge^k_E$ for a
bounded linear $Q^\sharp$ with $\norm{Q^\sharp}=\norm{\ev-\xi}\le1+\norm\xi$, using $\norm\ev\le1$. The corresponding $Q$
satisfies $Q(u,\dots,u)(\wedge^k_Ez)=\ev(z)(u,\dots,u)-\xi(z)(u,\dots,u)=\Lambda^k_\pi u(\wedge^k_Ez)$ for
$z\in\otimes^k_\pi E$, by Lemma~\ref{ob:lem:evaluation} and since $\xi(z)$ has zero diagonal; so
$Q(u,\dots,u)=\Lambda^k_\pi u$, since $\wedge^k_E$ is surjective.

\emph{Step 3 (the class).} Conditions (2) and (3) are equivalent by Lemma~\ref{ob:lem:split-criteria}(3).
\end{proof}

If $D$ is dense in a normed $K$-space $D'$, the canonical maps $\otimes^k_\pi D\to\otimes^k_\pi D'$ and
$\Lambda^k_\pi D\to\Lambda^k_\pi D'$ are isometric with dense range; so, with $\widehat D$ regarded as a normed $K$-space,
$\widehat\otimes^k_\pi\widehat D=\widehat\otimes^k_\pi D$ and $\widehat\Lambda^k_\pi\widehat D=\widehat\Lambda^k_\pi D$.
The maps are contractive. By continuity $\omega_D$ extends to a contractive map $\omega'\colon D'^k\to\widehat\Lambda^k_\pi D$,
alternating because tuples with two equal entries in $D'^k$ are limits of such tuples in $D^k$; its linearization
(Proposition~\ref{prop:exterior}) is a contraction $\Lambda^k_\pi D'\to\widehat\Lambda^k_\pi D$ restricting to the
inclusion on $\Lambda^k_\pi D$, so the maps are isometric. Density follows from
$x_1\wedge\dots\wedge x_k-y_1\wedge\dots\wedge y_k=\sum_ay_1\wedge\dots\wedge y_{a-1}\wedge(x_a-y_a)\wedge x_{a+1}\wedge\dots\wedge x_k$.
The tensor case is the same.

\begin{theorem}[ways for the class to vanish]\label{ob:thm:mechanisms}
Let $E$ and $D$ be normed spaces and $k\ge1$.
\begin{enumerate}
\item \textup{(Source.)} If $(\mathrm S^k_E)$ holds, then $\ob_k(E,D)=0$. More generally, $\ob_k(E,D)=0$ if
$\ev|_{\cN^k_E}=\tau\circ w^{\otimes k}|_{\cN^k_E}$ for an $\mathrm S^k$-split $w\colon E\to G$ and a bounded linear
$\tau\colon\cN^k_G\to\Mult^k_0(\cL(E,D);\widehat\Lambda^k_\pi D)$. This applies when $w\colon E\to G$ is
$\mathrm S^k$-split and $\varepsilon\colon\cL(E,D)\to\cL(G,\widehat D)$ is bounded linear with $\varepsilon(u)\circ w=u$
for every $u$, with
$\tau(\zeta)(u_1,\dots,u_k)=\widehat\wedge^k_{\widehat D}\bigl(\varepsilon(u_1)\otimes\dots\otimes\varepsilon(u_k)\bigr)\zeta$.
\item \textup{(Parameter.)} If the sequence $\cM^k(\cL(E,D);\widehat\Lambda^k_\pi D)$ splits, and in particular if
$\cL(E,D)$ is polarizable in degree $k$, then $\ob_k(E,D)=0$. More generally, $\ob_k(E,D)=0$ if there are a normed space $P$ with $(\Gamma^k_P)$, a
bounded linear map $\psi\colon\cL(E,D)\to P$ and a bounded $k$-linear map
$\Phi\colon P^k\to\cL(\widehat\otimes^k_\pi E,\widehat\Lambda^k_\pi D)$ such that $\Phi(h,\dots,h)$ vanishes on $\cN^k_E$
for $h\in P$ and $\Phi(\psi u,\dots,\psi u)=\widehat\Lambda^k_\pi u\circ\widehat\wedge^k_E$ for $u\in\cL(E,D)$.
\item \textup{(Target.)} Let
$\alt_*\colon\Mult^k_0(\cL(E,D);\widehat\Lambda^k_\pi D)\to\Mult^k_0(\cL(E,D);\widehat\otimes^k_\pi D)$ be
postcomposition with $\alt^k_D$. Then
$(\alt_*\circ\ev|_{\cN^k_E})_*[\cS^k_E]=0$. If $(\mathrm R^k_D)$ holds, then $\ob_k(E,D)=0$.
\item[(3$'$)] \textup{(Target, relative.)} If $w'\colon G\to\widehat D$ is $\mathrm R^k$-split and
$\varepsilon\colon\cL(E,D)\to\cL(E,G)$ is bounded linear with $w'\circ\varepsilon(u)=u$ for every $u$, then
$\ob_k(E,D)=0$.
\end{enumerate}
\end{theorem}
The proof is in Appendix~\ref{app:deferred}.

\begin{corollary}\label{ob:cor:mechanisms}
Each of $(\mathrm S^k_E)$, $(\mathrm R^k_D)$ and $(\Gamma^k_{\cL(E,D)})$ makes $(E,D)$ good for Banach values in degree $k$;
so does each of the data $(G,w,\varepsilon)$ of Theorem~\ref{ob:thm:mechanisms}(1) and $(G,w',\varepsilon)$ of
Theorem~\ref{ob:thm:mechanisms}(3$'$).
\end{corollary}

\begin{proof}
Theorems~\ref{ob:thm:mechanisms} and~\ref{ob:thm:obstruction}.
\end{proof}

\subsection{The ordering lift}\label{ob:ssec:ordering}

Throughout this subsection $E$ and $D$ are normed spaces, $k\ge1$, and $(e_t)_{t\in\Sigma}$ is a linearly independent
family in $E$, indexed by a totally ordered set $\Sigma$, whose span $V$ is dense in $E$; a Hamel basis of $E$ in any total
order is one. Write $x=\sum_tx(t)e_t$ for $x\in V$, $\Mult^k=\Mult^k(\cL(E,D);\hLambda^k_\pi D)$ and
$\Mult^k_0=\Mult^k_0(\cL(E,D);\hLambda^k_\pi D)$, and let $q$ and $\overline\ev$ be the maps defined after
Lemma~\ref{ob:lem:evaluation}. By the description of $\overline\ev$ given there, in the notation of the proof of
Theorem~\ref{ob:thm:obstruction}, a bounded linear $T\colon\hLambda^k_\pi E\to\Mult^k$ is $Q^\sharp$ for a lift $Q$ as in
Lemma~\ref{fib:lem:test} if and only if $q\circ T=\overline\ev$. Let $c_k(E,D)\in[0,\infty]$ be the infimum of $\norm Q$
over these lifts; by Lemma~\ref{fib:lem:test} it is finite if and only if $(E,D)$ is good for Banach values in degree $k$.

Let $\Sigma^k_<\subseteq\Sigma^k$ be the set of strictly increasing $t$, and put $e^\wedge_t=e_{t_1}\wedge\dots\wedge e_{t_k}$
and $e^\otimes_t=e_{t_1}\otimes\dots\otimes e_{t_k}$ for $t\in\Sigma^k$; the $e^\wedge_t$ with $t\in\Sigma^k_<$ form a basis
of $\Lambda^k_KV$. Let $\wedge_V\colon\otimes^k_KV\to\Lambda^k_KV$ be the quotient map and
$j_\Lambda\colon\Lambda^k_KV\to\hLambda^k_\pi E$, $j_\otimes\colon\otimes^k_KV\to\widehat\otimes^k_\pi E$ the canonical maps,
so that $\widehat\wedge^k_E\circ j_\otimes=j_\Lambda\circ\wedge_V$, and put $\norm y_\pi=\norm{j_\Lambda y}$. By the paragraph
before Theorem~\ref{ob:thm:mechanisms}, $\norm y_\pi$ is the projective seminorm of $y$ computed from representations with
entries in $V$, and $j_\Lambda(\Lambda^k_KV)$ is dense in $\hLambda^k_\pi E$; the same holds for $j_\otimes$. Norms in
$[0,\infty]$ of linear maps on $\Lambda^k_KV$, for $\norm\cdot_\pi$, and of $k$-linear maps on $V^k$ are best constants,
as are the norms of lifts $Q$ and of the maps $Q^\sharp$.

The \emph{ordering section} is the linear map $s_<\colon\Lambda^k_KV\to\otimes^k_KV$ with $s_<(e^\wedge_t)=e^\otimes_t$ for
$t\in\Sigma^k_<$. Then $\wedge_V\circ s_<=\id$, and $s_<=\cP^{(k)}\circ\alt^k$ for the map $\cP^{(k)}$ of
Proposition~\ref{spl:prop:ordering} and the antisymmetrization $\alt^k\colon\Lambda^k_KV\to\otimes^k_KV$. The
\emph{ordering lift} assigns to $u_1,\dots,u_k\in\cL(E,D)$ the linear map
$Q_<(u_1,\dots,u_k)\colon\Lambda^k_KV\to\Lambda^k_\pi D$ with
\begin{equation}\label{ob:eq:ordering-lift}
 Q_<(u_1,\dots,u_k)(e^\wedge_t)=u_1e_{t_1}\wedge\dots\wedge u_ke_{t_k}\qquad(t\in\Sigma^k_<);
\end{equation}
on $x_1\wedge\dots\wedge x_k$ with $x_1,\dots,x_k\in V$ it is the right side of~\eqref{ord:eq:canonical-lift}. Let
$Q^\sharp_<\colon\Lambda^k_KV\to\Mult^k$ be the linear map with $Q^\sharp_<(y)(u_1,\dots,u_k)=Q_<(u_1,\dots,u_k)y$; it takes
values in $\Mult^k$ by~\eqref{ob:eq:ordering-lift}. A \emph{zero-diagonal correction} is a linear map
$L\colon\Lambda^k_KV\to\Mult^k_0$. The proofs of Theorems~\ref{ob:thm:ordering} and~\ref{ob:thm:ultrametric},
Lemma~\ref{ob:lem:repeated} and Corollaries~\ref{ob:cor:adjacent} and~\ref{ob:cor:z-linear} are in
Appendix~\ref{app:deferred}.

\begin{theorem}[ordering representation]\label{ob:thm:ordering}
\begin{enumerate}
\item $\ev\circ j_\otimes\circ s_<=Q^\sharp_<$ and $q\circ Q^\sharp_<=\overline\ev\circ j_\Lambda$.
\item For every lift $Q\colon\cL(E,D)^k\to\cL(\Lambda^k_\pi E,\hLambda^k_\pi D)$ as in Lemma~\ref{fib:lem:test},
$L=Q^\sharp\circ j_\Lambda-Q^\sharp_<$ is a zero-diagonal correction; equivalently,
$Q(u_1,\dots,u_k)\circ j_\Lambda=Q_<(u_1,\dots,u_k)+\Phi(u_1,\dots,u_k)$, where $\Phi(u_1,\dots,u_k)y=L(y)(u_1,\dots,u_k)$ is
$k$-linear in $(u_1,\dots,u_k)$ and $\Phi(u,\dots,u)=0$ for $u\in\cL(E,D)$.
\item For every zero-diagonal correction $L$ with $\norm{Q^\sharp_<+L}<\infty$ there is a unique lift $Q$ with
$Q^\sharp\circ j_\Lambda=Q^\sharp_<+L$, and $\norm Q=\norm{Q^\sharp_<+L}$. Hence $c_k(E,D)=\inf_L\norm{Q^\sharp_<+L}$ over the
zero-diagonal corrections $L$, and $(E,D)$ is good for Banach values in degree $k$ if and only if $\norm{Q^\sharp_<+L}<\infty$
for some zero-diagonal correction $L$.
\item \textup{(Defect form.)} The formula
$\Omega_<(x)=\ev(x_1\otimes\dots\otimes x_k)-Q^\sharp_<(x_1\wedge\dots\wedge x_k)$, for $x\in V^k$, defines a $k$-linear map
$\Omega_<\colon V^k\to\Mult^k_0$. For every zero-diagonal correction $L$, with $(L\circ\wedge)(x)=L(x_1\wedge\dots\wedge x_k)$,
\[
 \norm{Q^\sharp_<+L}\le1+\norm{\Omega_<-L\circ\wedge},\qquad\norm{\Omega_<-L\circ\wedge}\le1+\norm{Q^\sharp_<+L}.
\]
So $(E,D)$ is good for Banach values in degree $k$ if and only if $\Omega_<$ is the sum of a bounded $k$-linear map
$V^k\to\Mult^k_0$ and an alternating $k$-linear map $V^k\to\Mult^k_0$. The infima $c_k(E,D)$ and
$\inf_L\norm{\Omega_<-L\circ\wedge}$ are finite together, and then they differ by at most~$1$.
\end{enumerate}
\end{theorem}

If $\norm{\cP_<(x\otimes y)}\le C\norm x\norm y$ for $x,y\in V$, as in Proposition~\ref{spl:prop:ordering}, no correction is
needed: the lift $Q$ of
Proposition~\ref{fib:prop:split}(1) built from the section $s=\cP^{(k)}\circ\alt^k_E$ of
Proposition~\ref{spl:prop:ordering}(2) satisfies $Q(u_1,\dots,u_k)\circ j_\Lambda=Q_<(u_1,\dots,u_k)$, since
$s\circ j_\Lambda=j_\otimes\circ s_<$.

\begin{lemma}\label{ob:lem:repeated}
The space $\cN^k_E$ is the closure in $\widehat\otimes^k_\pi E$ of the span of the pure tensors $x_1\otimes\dots\otimes x_k$
with $x_a=x_{a+1}$ for some $a<k$, and also of the span of the pure tensors with two equal entries. The pure tensors of the
first kind with all entries in $V$ span a dense subspace of $\cN^k_E$.
\end{lemma}

\begin{corollary}[adjacent diagonals]\label{ob:cor:adjacent}
The pair $(E,D)$ is good for Banach values in degree $k$ if and only if there is a bounded $k$-linear
$\Psi\colon E^k\to\Mult^k_0$ with
\[
 \Psi(x_1,\dots,x_k)(u_1,\dots,u_k)=u_1x_1\wedge\dots\wedge u_kx_k
\]
whenever two of the $x_a$ coincide, equivalently whenever $x_a=x_{a+1}$ for some $a<k$. A lift $Q$ gives such a $\Psi$ with
$\norm\Psi\le1+\norm Q$, and such a $\Psi$ gives a lift $Q$ with $\norm Q\le1+\norm\Psi$.
\end{corollary}

\begin{corollary}[z-linear form]\label{ob:cor:z-linear}
Let $\delta>0$, and let $\nu\colon\hLambda^k_\pi E\to\widehat\otimes^k_\pi E$ be a homogeneous map, $\nu(cy)=c\,\nu(y)$ for $c\in K$,
with $\widehat\wedge^k_E\circ\nu=\id$ and $\norm{\nu(y)}\le(1+\delta)\norm y$; it exists by Lemma~\ref{ob:lem:metric-surjection},
choosing $\nu(y)$ for one point $y$ of each orbit $K^\times y$, $y\ne0$, and extending homogeneously. Put
$\Omega^\nu_<=\ev\circ\nu\circ j_\Lambda-Q^\sharp_<\colon\Lambda^k_KV\to\Mult^k$. Then $\Omega^\nu_<$ takes values in $\Mult^k_0$, is
homogeneous, and is \emph{z-linear}:
\[
 \Bigl\lVert\Omega^\nu_<\Bigl(\sum_iy_i\Bigr)-\sum_i\Omega^\nu_<(y_i)\Bigr\rVert\le2(1+\delta)\sum_i\norm{y_i}_\pi
\]
for every finite family $(y_i)$ in $\Lambda^k_KV$. Let $\operatorname{dist}(\Omega^\nu_<)$ be the infimum, over the linear maps
$L\colon\Lambda^k_KV\to\Mult^k_0$, of the infimum of the $C$ with $\norm{\Omega^\nu_<(y)-L(y)}\le C\norm y_\pi$ for all $y$. Then
$\operatorname{dist}(\Omega^\nu_<)$ and $c_k(E,D)$ are finite together, and then they differ by at most $1+\delta$. In
particular $\ob_k(E,D)=0$ if and only if $\Omega^\nu_<$ lies at finite distance from a linear map $\Lambda^k_KV\to\Mult^k_0$.
\end{corollary}

\begin{theorem}[nonarchimedean values]\label{ob:thm:ultrametric}
Let $K$ be nonarchimedean, and suppose that $C_e=\sup_t\norm{e_t}<\infty$ and that $\abs{x(t)}\le C_c\norm x$ for $x\in V$,
$t\in\Sigma$ and a constant $C_c$. Let $Y$ be a Banach space with a nonarchimedean norm $\lvert\!\lvert\!\lvert\cdot\rvert\!\rvert\!\rvert_Y$ and constants $0<c_-\le c_+$
with $c_-\lvert\!\lvert\!\lvert\cdot\rvert\!\rvert\!\rvert_Y\le\norm\cdot\le c_+\lvert\!\lvert\!\lvert\cdot\rvert\!\rvert\!\rvert_Y$, and let $\alpha\colon\cN^k_E\to Y$ be bounded linear. Then there is a bounded linear
$\tilde\alpha\colon\widehat\otimes^k_\pi E\to Y$ of norm at most $2(c_+/c_-)\norm\alpha(C_cC_e)^k$ with
$\tilde\alpha(j_\otimes z)=\alpha\bigl(j_\otimes(z-s_<(\wedge_Vz))\bigr)$ for $z\in\otimes^k_KV$ and $\tilde\alpha=\alpha$ on $\cN^k_E$. Hence
$\alpha_*[\cS^k_E]=0$. In particular $\ob_k(E,D)=0$ whenever $\hLambda^k_\pi D$ has an equivalent nonarchimedean norm.
\end{theorem}
